\subsection{余积}

在本节中，X 表示一个 \(C^{r}\) 类流形，x 是 X 的一个点。

13.5.1. 设 \(k \leqslant r\)。若 \(\Delta\) 是对角映射 \(y \mapsto (y, y)\)，即从 X 到 \(X \times X\) 的对角映射，则 \(\Delta_*\) 把 \(\mathrm{T}_x^{(k)}(\mathrm{X})\) 映入 \(\mathrm{T}_{(x,x)}^{(k)}(\mathrm{X} \times \mathrm{X})\)，后者是 \(\mathrm{T}_{x}^{(k)}(\mathrm{X})\otimes\mathrm{T}_{x}^{(k)}(\mathrm{X})\) 的一个向量子空间（参见 13.4.5）。这样我们就得到一个线性映射，仍记作 \(\Delta_{*}\)（或 \(c\)）：

\[
\mathrm{T} _ {x} ^ {(k)} (\mathrm{X}) \rightarrow \mathrm{T} _ {x} ^ {(k)} (\mathrm{X}) \otimes \mathrm{T} _ {x} ^ {(k)} (\mathrm{X});
\]

我们把它称为附加于 \(\mathbf{T}_{x}^{(k)}(\mathbf{X})\) 的余积。装备了这个余积之后，\(\mathbf{T}_{x}^{(k)}(\mathbf{X})\) 是一个余结合且余交换的余代数（A, III, p. 144–145）；它以把每个点分布映为其常数项的线性形式（13.2.1）作为余单位。若 \(k' \leqslant k\)，则 \(\mathbf{T}_{x}^{(k')}(\mathbf{X})\) 到 \(\mathbf{T}_{x}^{(k)}(\mathbf{X})\) 的包含映射是余代数的态射。

若 \(c = (\mathrm{U},\varphi ,\mathrm{E})\) 是 \(\mathbf{X}\) 的以 \(x\) 为中心的卡，则同构

\[
\theta_ {c} \colon \mathrm{T} ^ {(k)} (\mathrm{E}) \rightarrow \mathrm{T} _ {x} ^ {(k)} (\mathrm{X})
\]

是一个余代数同构，其中 \(\mathrm{T}^{(k)}(\mathrm{E})\) 装备的是由 TS(E) 的余积诱导的余积（参见 A, IV, § 5, no. 7）。

若 \(\varphi: X \to Y\) 是一个 \(C^r\) 类流形态射，则从 \(T_x^{(k)}(X)\) 到 \(T_{\varphi(x)}^{(k)}(Y)\) 的、由 \(\varphi_*\) 诱导的映射是余代数的态射。

13.5.2. 设 \(F_{1}\)、\(F_{2}\) 与 F 是分离的多范数空间，并设 \((u_{1}, u_{2}) \mapsto u_{1}.u_{2}\) 是从 \(F_{1} \times F_{2}\) 到 F 的一个连续双线性映射。设 \(t \in \mathrm{T}_{x}^{(r)}(\mathrm{X})\)，并设它在余积下的像为 \(\boldsymbol{c}(t) = \sum_{j} u_{j} \otimes v_{j}\)，其中 \(u_{j},v_{j} \in \mathrm{T}_{x}^{(r)}(\mathrm{X})\)。设 \(f_{i}: X \to F_{i}\)（\(i=1,2\)）是 \(C^{r}\) 类映射。我们有

\[
\langle f _ {1}. f _ {2}, t \rangle = \sum_ {j} \langle f _ {1}, u _ {j} \rangle . \langle f _ {2}, v _ {j} \rangle ,
\]

按记号的滥用，我们把它写成：

\[
\langle f _ {1}. f _ {2}, t \rangle = \langle f _ {1} \otimes f _ {2}, c (t) \rangle .
\]

13.5.3. 设 \(t \in \mathbf{T}_x^{(k)}(\mathbf{X})\)，其中 \(k \leqslant r\)。

\begin{enumerate}
\def\labelenumi{\alph{enumi})}
\item
  要使 \(c(t)=t\otimes t\)，必须且只需 \(t\) 等于 0 或等于 \(\varepsilon_{x}\)。
\item
  要使 \(c(t)=t\otimes\varepsilon_{x}+\varepsilon_{x}\otimes t\)，必须且只需 \(t\) 是一个切向量，即 \(t\in\mathrm{T}_{x}^{(1)+}(X)\)。
\end{enumerate}

